% Unit 1 map -- one page: calendar, readings, objectives.
\documentclass{apchem-worksheet}

\apcunit{1}
\apcunittitle{Atomic Structure and Properties}
\apcdoctitle{Unit Map}
\apcsubtitle{CED Unit 1 \textbullet\ 7--9\% of the AP Exam \textbullet\ 5 blocks + exam}

\begin{document}
\apctitleblock

\begin{apcnote}[How this unit is graded]
Unit exam Monday, Aug 24: 20 multiple choice (\textbf{no calculator}) + 2 short
free-response, 48 minutes. Worksheets are practice, not grades --- but every
exam question traces back to one.
\end{apcnote}

\section*{\sffamily Block calendar}

\renewcommand{\arraystretch}{1.35}
\noindent\begin{tabular}{@{}p{2.1cm}p{5.1cm}p{1.6cm}p{3.4cm}p{2.2cm}@{}}
\toprule
\textbf{Date} & \textbf{Topics} & \textbf{CED} & \textbf{Read (PDF pp.)} & \textbf{Practice}\\
\midrule
Wed Aug 12 & Moles, molar mass; mass spectra of elements & 1.1, 1.2 & \S3.2--3.5 \textbullet\ 112--123 & WS 1\\
Fri Aug 14 & Percent composition; empirical \& molecular formulas; mixtures & 1.3, 1.4 & \S3.6--3.7 \textbullet\ 123--132; \S1.10 \textbullet\ 53--57 & WS 1--2\\
Mon Aug 17 & Atomic structure; electron configurations; Coulomb's law & 1.5 & \S7.5, 7.8--7.11 \textbullet\ 346--364 & WS 3\\
Wed Aug 19 & Photoelectron spectroscopy & 1.6 & \S7.12 \textbullet\ 368--369 & WS 3\\
Fri Aug 21 & Periodic trends; valence electrons \& ionic compounds & 1.7, 1.8 & \S7.12 \textbullet\ 364--371; \S8.4 \textbullet\ 400--404 & WS 4 (FRQ)\\
\textbf{Mon Aug 24} & \textbf{UNIT 1 EXAM} & all & review sheet & ---\\
\bottomrule
\end{tabular}

\vspace{0.6em}
\section*{\sffamily By the end of this unit you can\ldots}

\begin{itemize}[leftmargin=*,itemsep=0.15em]
  \item Convert among grams, moles, and particles, and justify each step
        with dimensional analysis. \ced{1.1}
  \item Read a mass spectrum: estimate average atomic mass from peak
        positions and heights, and identify the element. \ced{1.2}
  \item Determine empirical and molecular formulas from percent composition
        or combustion-style data. \ced{1.3}
  \item Distinguish pure substances from mixtures using particulate
        diagrams, and compute composition of mixtures. \ced{1.4}
  \item Write ground-state electron configurations for atoms and ions, and
        spot excited states. \ced{1.5}
  \item Interpret a photoelectron spectrum: peak position $\to$ binding
        energy, peak area $\to$ electron count; identify the element. \ced{1.6}
  \item Explain radius, ionization energy, and electronegativity trends
        using Coulomb's law --- not just recite them. \ced{1.7}
  \item Predict ionic charges from electron configurations and write
        correct ionic formulas. \ced{1.8}
\end{itemize}

\begin{apctrap}[Where students lose points in this unit]
Averaging isotope masses without weighting; confusing peak \emph{height}
(abundance) with peak \emph{position} (mass) on spectra; writing configurations
for ions as if they were atoms; ranking trends correctly but ``explaining''
them by restating the trend instead of using nuclear charge, distance, and
shielding.
\end{apctrap}

\end{document}
